\section{第 4 章 系统总体设计}

\subsection{4.1 系统总体架构}

\begin{figure}[htbp]
\centering
\caption{图 4-1  智能气象信息服务系统总体架构图}
\label{fig:4-1}
\end{figure}

\begin{figure}[htbp]
\centering
\caption{图 4-2  系统核心数据流图}
\label{fig:4-2}
\end{figure}

\begin{figure}[htbp]
\centering
\caption{图 4-3  系统主要模块关系图}
\label{fig:4-3}
\end{figure}

\begin{figure}[htbp]
\centering
\caption{图 4-4  首页导航页界面示意图1}
\label{fig:4-4}
\end{figure}

\begin{figure}[htbp]
\centering
\caption{图 4-5  首页导航页界面示意图2}
\label{fig:4-5}
\end{figure}

本系统使用基于Flask的Web应用架构，整体可以分为路由层、服务层、分析层、模型层、展示层五个部分。路由层主要是处理页面入口以及接口调度，服务层主要是对外部数据进行抓取、整合以及运行时的配置，分析层主要是对历史气候、探空参数、农业指标等进行计算逻辑的处理，模型层主要是机器学习误差校正模型的加载和使用，展示层主要是由HTML模板、前端脚本、图表组件组成，实现数据可视化以及交互展示。

app.py 是系统主入口，主要用来组织首页、多模式预报、探空分析、历史气候、农业气象等各个页面的路由，并对相关的接口进行统一管理。以该入口文件为基准，系统把不同的业务逻辑拆分成独立的模块，多模式预报用advanced\_forecast\_service.py来处理，历史气候分析用history\_analyzer.py来处理，探空分析用sounding\_parser.py、sounding\_analyzer.py和sounding\_plotter.py协同处理，机器学习误差校正用ml\_correction.py和train\_bias\_correction.py支持，农业气象能力由作物知识库和相关计算模块提供。

\subsection{4.2 模块划分与职责}

系统是由五个主要业务模块构成的。

第一部分是多模式预报模块。该模块从Open-Meteo等处获取多种数值预报模式的结果，对逐小时预报进行组织、统一字段映射、预报时次估计以及缓存处理，最后向前端提供可以直接展示的数据结构。

第二，探空分析模块。从怀俄明大学探空接口抓取指定站点、指定时次的探空资料，解析层结数据，判断是否有效观测，计算稳定度参数，产生图形输出。

第三部分是历史气候分析模块。该模块主要是对历史站点数据进行解析、清洗、聚合，完成年度极值统计、长期趋势分析、EW指数计算、代表事件筛选等工作。

第四，机器学习误差校正模块。本模块完成训练样本准备、观测真值对齐、特征提取、模型训练、模型加载和运行时校正的应用，是提高预报可信度的关键部分。

第五章农业气象服务模块。该模块在天气及历史分析结果的基础上，采用作物知识库、积温计算模型等手段，完成农业风险预警、农事建议的输出。

除上述五个核心业务模块外，系统还专门设置了首页导航页作为平台级入口页面。该页面并不承担复杂分析计算，但负责整合实时天气概览、统一导航、智能提醒和各主要模块跳转关系，是用户理解系统整体功能结构和进入后续业务流程的起点。因此，在系统设计中，首页不是可有可无的装饰性页面，而是连接各模块的重要枢纽。

\subsection{4.3 数据流设计}

系统数据流可以简要归纳为外部数据获取、本地清洗和组织、分析计算、接口封装、页面展示这五个过程。

在多模式预报流程当中，系统会先去开放预报接口申请 ECMWF，GFS，ICON 等模式的数据，之后依照统一的目标站点坐标和字段组织准则来创建标准化的输出。对于ECMWF数据，如果模型已经加载并且可以识别出预报起报时次，则系统会调用机器学习校正逻辑，得到校正后的结果，最后将原始结果和校正结果一起返回给前端。

在探空分析流程中，系统先根据页面输入的北京时间换算为标准 UTC 探空时次，再向怀俄明大学接口请求数据。若指定时次无有效层结数据，则允许回退到前一个标准时次；若解析成功，则进一步提取层结表、计算稳定度指数并生成 Skew-T 或其他图形。

在历史气候流程中，系统从本地历史观测文件中读取多年样本，完成日期解析、变量清洗和统计聚合，再生成年度统计、趋势结果和 EW 指数代表事件结构，以供年度页和趋势页展示。

在农业气象流程中，系统则将预报结果与作物知识库结合，计算积温进度、成熟时间估计及风险提示信息，并将结果展示在农业页面中。

\subsection{4.4 站点与时次统一策略}

由于系统同时涉及预报数据、观测数据、历史数据和探空数据，若站点配置和时间口径不统一，容易造成结果解释混乱。因此，系统在设计中强调统一的站点和时次策略。

在站点方面，系统将杭州国家站作为主要目标站点，并在预报服务与机器学习误差校正链路中尽量统一到相同的经纬度配置，以避免出现“训练一个点、运行另一个点”的情况。

时间上将本地时间、UTC 标准时次、预报起报时间、目标预报时间加以区别。探空分析页面输入以北京时间为准，但是后端会转换成UTC并统一格式为00Z或者12Z；机器学习误差校正部分使用issue time和target time进行明确的区分，防止训练集和推理阶段口径不同。

该种统一策略一方面可以提高系统的实现一致性，另一方面也可以给论文的结果解释提供更加可靠的基础。

